\documentclass[10pt]{article}
\usepackage[a4paper,margin=2.2cm]{geometry}
\usepackage{amsmath,amssymb,bm,booktabs,array,xcolor}
\usepackage[hidelinks]{hyperref}
\usepackage{natbib}

\newcommand{\vc}[1]{\bm{#1}}
\newcommand{\mt}[1]{\mathbf{#1}}
\newcommand{\sk}[1]{[#1]_{\times}}

\title{Phase 2: Simulator and verification\\[2pt]
\large Fault tolerant control and post fault trajectory planning of a redundant cable driven parallel robot}
\author{M.\,J.\ Tachia, A.\,V.\ Maloletov}
\date{Technical report, 26 October 2026}

\begin{document}
\maketitle

\noindent\fbox{\parbox{0.97\linewidth}{\small\textit{About this document.} This document describes the simulator on which every result of the study will be produced and reports the verification gates of the specification. It builds on the Phase~1 methods (\texttt{docs/phase1/methods.pdf}), whose equation numbers are referred to as (M\,$n$). Two changes to the Phase~1 model came out of this phase and are already carried into the methods: the orientation is integrated as a unit quaternion, and the damping of a cable fades to zero at zero stretch. One departure from the specification is made on evidence: the plant step is 62.5~$\mu$s, below the range of 0.1 to 0.5~ms written in the specification, because the step convergence test fails above it once a cable breaks. The closed loop runs here use a verification controller, a motor angle PD with tension feedforward, which is not the method of the study; its only purpose is to exercise the simulator before Phase~3. All results carry the code digest \texttt{c77663b917040fde}.}}

% =====================================================================
\section{Simulation model}\label{sec:model}

The simulator integrates the full model of Phase~1, the cable law (M\,13) and the winch equation (M\,17): a rigid platform with six degrees of freedom, eight elastic cables that can only pull, and eight geared winches whose motors are driven by current controlled drives. Three additions turn that model into a simulator that can be trusted through a fault, and a fourth defines what the controller is allowed to see.

\paragraph{Orientation.} The plant state is
\begin{equation}\label{eq:state}
\vc x=\big[\vc p^\top,\ \vc\eta^\top,\ \vc v^\top,\ \vc\omega^\top,\ \vc\theta^\top,\ \dot{\vc\theta}^\top\big]^\top\in\mathbb R^{13+2n},
\qquad
\dot{\vc\eta}=\tfrac12\begin{bmatrix}-\vc\omega^\top\vc\epsilon\\ \eta_0\vc\omega+\vc\omega\times\vc\epsilon\end{bmatrix},
\end{equation}
where $\vc\eta=[\eta_0,\vc\epsilon^\top]^\top$ is the unit quaternion of the platform and $\vc\omega$ its angular velocity in the world frame. The quaternion is renormalised after every integration step. The Euler angles of Phase~1 have a singularity at a pitch of $\pm90^\circ$, which a controlled platform never approaches but an uncontrolled one can reach after an upper cable breaks, as the runs of Section~\ref{sec:convergence} show; with the quaternion the plant stays regular through any rotation, and the ZYX angles are recovered from $\vc\eta$ for the estimators, the controller and the logs. Everything else in the platform dynamics is (M\,12) unchanged.

\paragraph{Cables near slack.} A cable becomes slack whenever its stretch reaches zero, and after a fault this happens many times, as cables lose their load and catch it again. In the plain Kelvin--Voigt law a cable that tightens at the stretch rate $\dot\Delta$ jumps at once to the force $\beta k\dot\Delta$, which is neither physical nor harmless for a fixed step integrator, since a jump inside a step costs the integrator its order. The tension law (M\,13) therefore carries the factor $s(\Delta)=3r^2-2r^3$, $r=\min(\Delta/\Delta_s,1)$, which lets the damping grow from zero with the stretch, as Hunt and Crossley proposed for impacts \citep{hunt1975coefficient}. With $\Delta_s=0.1$~mm the factor equals one above a tension of 0.4 to 2.2~N, depending on the cable length, so every taut state of the robot follows the identified law. The effect on the solution is not small: in the stiffest scenario of Section~\ref{sec:convergence}, the peak tension after a lower cable breaks is 244~N with the factor and 360~N without it, everything else being equal, so the extra 116~N comes from the jump alone.

\paragraph{Drive limits.} The drive applies the commanded torque $\tau_i$ only up to the peak torque of the motor, and while the motor is driving ($\tau_i\dot\theta_i>0$) it lets the torque fall linearly to zero over the last 10\% below the maximum speed,
\begin{equation}\label{eq:drive}
\tau_i^{\mathrm{drv}}=\begin{cases}\bar\tau_i\,\min\!\Big(1,\dfrac{\dot\theta_{\max}-|\dot\theta_i|}{0.1\,\dot\theta_{\max}}\Big)^{+}, & \bar\tau_i\dot\theta_i>0,\\[4pt] \bar\tau_i, & \text{otherwise,}\end{cases}
\qquad \bar\tau_i=\operatorname{sat}_{\tau_{\mathrm{pk}}}(\tau_i),
\end{equation}
with $\tau_{\mathrm{pk}}=4.0$~N\,m and $\dot\theta_{\max}=9000$~rpm for the MSK030C \citep{rexroth_msk2018}. The faults (M\,18) act after the drive, $\tau_i^{\mathrm{act}}=(1-\rho_i)\tau_i^{\mathrm{drv}}+f_i$, so the drive's own current, and with it the measured torque, stays $\tau_i^{\mathrm{drv}}$ whatever happens downstream. The speed limit matters after a break: the winch of the broken cable loses its load, and a tension loop that keeps asking for its reference would otherwise spin the motor without bound. The shape of the limiter is our choice, since the limiter of the drive is not documented.

\paragraph{Drum windows.} The drum holds 6~m of cable, which the methods turned into a length window $l_i^-\le l_{u,i}\le l_i^-+\Delta l_{\max}$ for each cable (M\,15). The plant does not stop at the window but flags every sample at which a drum would be overfull or empty, so that a run that leaves the windows is recorded as such.

\paragraph{Sensors and timing.} Each winch reports two signals, the motor angle from its multi turn absolute encoder and the motor current, and no cable carries a load cell. Table~\ref{tab:sensors} lists the sensor model. The encoder option fitted to IPAnema~2 is not published; we take the M1 option of the MSK motors, which is multi turn and absolute as \citet{pott2013ipanema} state, and its resolution of $128\times2^{13}=2^{20}$ counts per revolution, the figure that the manufacturer's own example gives for this option \citep{rexroth_msk2018}. Through the gearbox one count is 25~nm of cable, so the quantisation is far below anything the platform could resolve, but it is not irrelevant to the observer of each winch, whose third channel differentiates the angle three times. The system accuracy of $\pm120''$ enters as an error that repeats once per motor revolution with a random phase per winch,
\begin{equation}\label{eq:encoder}
\tilde\theta_i=q_e\Big\lfloor\frac{\theta_i+a_e\sin(\theta_i+\varphi_i)}{q_e}\Big\rfloor,\qquad q_e=\frac{2\pi}{2^{20}},\quad a_e=120'' ,
\end{equation}
and the measured torque is $\tilde\tau_i=\tau_i^{\mathrm{drv}}+\nu_i$ with white Gaussian noise $\nu_i$. No noise figure has been published for the current measurement of the drive, so its standard deviation, 0.004~N\,m or 0.96~N of cable tension, is a design choice that the evaluation doubles, as the specification asks.

\begin{table}[t]
\centering
\caption{Sensor model and loop timing. The resolution is the manufacturer's example for the M1 option; the current noise is a design choice evaluated at one and two times its value.}\label{tab:sensors}
\small
\begin{tabular}{@{}>{\raggedright\arraybackslash}p{0.17\linewidth}>{\raggedright\arraybackslash}p{0.4\linewidth}>{\raggedright\arraybackslash}p{0.37\linewidth}@{}}
\toprule
Quantity & Value & Source or status\\
\midrule
Encoder & M1: Hiperface, multi turn (4096 rev.), 128 signal periods & \citet{rexroth_msk2018}, Sec.~4.3\\
Resolution & $2^{20}$ counts/rev (25~nm of cable) & manufacturer's example; depends on the drive\\
Accuracy & $\pm120''$, once per revolution, random phase & \citet{rexroth_msk2018}; error shape our choice\\
Torque measurement & $\tau^{\mathrm{drv}}$ + noise, $\sigma=0.004$~N\,m (0.96~N) & design choice; evaluated at $1\times$ and $2\times$\\
Delay & one inner sample (1~ms) & specification\\
Inner rate & 1~kHz, zero order hold & 2~ms set points on IPAnema \citep{pott2013ipanema}, 1~ms on IPAnema~3 \citep{kraus2015hybrid}\\
Plant step & 62.5~$\mu$s, RK4, 16 steps per sample & step convergence test (Section~\ref{sec:convergence})\\
\bottomrule
\end{tabular}
\end{table}

At every inner sample $t_k=k\cdot1$~ms the simulator first switches on any fault whose onset has been reached, then reads the encoders and the drive torque from the true state, hands the controller the reading taken one sample earlier, receives the torque command, and integrates the plant to $t_{k+1}$ with the command held, in 16 RK4 steps. Faults therefore start on the millisecond grid, which is the resolution at which the controller could notice them in any case.

% =====================================================================
\section{Implementations}\label{sec:impl}

The plant exists twice. The NumPy implementation follows the equations line by line and is the reference that the tests check. The compiled implementation, written with CasADi \citep{andersson2019casadi}, computes the same quantities and is the one that runs the campaign. It solves the twist equations in closed form instead of factorising $\mt M_\nu$: subtracting $\sk{\vc c_w}$ times the first block row of (M\,12) from the second gives
\begin{equation}\label{eq:closedform}
\mt I_c\,\dot{\vc\omega}=\vc r_\omega-\vc c_w\times\vc r_v,\qquad \dot{\vc v}=\vc r_v/m+\vc c_w\times\dot{\vc\omega},
\end{equation}
where $\vc r_v$ and $\vc r_\omega$ are the force and moment on the right of (M\,12) minus the bias terms, and $\mt I_c^{-1}=\mt R\mt I_G^{-1}\mt R^\top$ is available without inversion at run time. One call advances the state over a whole millisecond, so Python is entered once per inner sample. Over 200 random states with large rotations, slack and broken cables, saturated drives and winches near their speed limit, the two implementations agree to a relative difference of $5\times10^{-13}$ in $\dot{\vc x}$, and over a closed loop run without sensor noise their trajectories agree to round off. The compiled plant runs a 16~s closed loop run on the reference T1 in 11.7~s; the NumPy plant, at the 1.7~ms it needs per RK4 step, would take about seven minutes.

The estimators and the MPC need a third model, the rigid cable platform of (M\,12) with the tensions as inputs. It is built in CasADi as well, discretised by one RK4 step, and returns its Jacobians with respect to the state and the tensions by algorithmic differentiation; the tests compare them with central differences and the dynamics with the NumPy rigid body model. The same module provides the unstretched lengths $EA\,l/(EA+t)$ that the encoders measure under a given tension, which is the measurement function of the IMM (M\,29).

% =====================================================================
\section{Verification}\label{sec:verification}

Table~\ref{tab:gates} collects the gates of the specification and adds the ones this phase introduced. All of them pass on the study parameters, with one qualification in the step convergence test that the following subsection explains.

\begin{table}[t]
\centering
\caption{Verification gates on the study parameters. Gates marked S are those of the specification; the others were added in this phase.}\label{tab:gates}
\begin{tabular}{@{}>{\raggedright\arraybackslash}p{0.43\linewidth}>{\raggedright\arraybackslash}p{0.2\linewidth}>{\raggedright\arraybackslash}p{0.3\linewidth}@{}}
\toprule
Check & Criterion & Result\\
\midrule
S\quad Length Jacobian against central differences & rel.\ error $<10^{-7}$ & $9.0\times10^{-10}$ (Phase~1)\\
S\quad Energy, locked winches, undamped cables, no input, 10~s at 62.5~$\mu$s & drift $<10^{-6}$ of initial kinetic energy & $3.5\times10^{-11}$\\
S\quad $\dot{\mt M}-2\mt C$ skew symmetric, $\dot{\mt M}$ exact & $<10^{-10}$ & $7.1\times10^{-15}$\\
S\quad Static equilibrium at home, study parameters & matches the LP & tensions 151.4 / 58.6~N, $\max|\dot{\vc x}|=4.8\times10^{-10}$\\
S\quad Step halving & metrics change $<0.1\%$ & passes at 62.5~$\mu$s; one exception (Section~\ref{sec:convergence})\\
S\quad Reference solution (Radau IIA, rtol $10^{-9}$) & & $5.8\times10^{-8}$~m, 0.010~N after a break\\
S\quad MuJoCo, same inputs, 5~s & difference $<1\%$ of the motion & $1.4\times10^{-10}$ of the motion, with slack cables\\
\quad Compiled plant against NumPy plant, 200 states & & $5.0\times10^{-13}$ relative\\
\quad Prediction model Jacobians against central differences & abs.\ error $<10^{-7}$ & pass\\
\quad Sensor model, loop delay, fault onset, drive limits, drum flags & exact values & pass\\
\bottomrule
\end{tabular}
\end{table}

\subsection{Step convergence}\label{sec:convergence}

The convergence test runs the closed loop on T1 for 8~s with the sensors on and repeats it with the plant step halved from 1~ms down to 31.25~$\mu$s. Five scenarios cover the range of what the campaign will ask of the integrator: no fault; a loss of effectiveness of 60\% on winch~3 at 4~s; a break of the lower cable~5 at 4~s; a break of the upper cable~1 at 4~s; and a break of cable~5 with twice the nominal stiffness and the larger retardation time of 0.03~s, which are the fastest cable dynamics in the robustness ranges. The two most violent cases are also run without sensor noise and quantisation, so that the integrator can be judged apart from the sensitivity of the loop to single encoder counts. The metrics compared on halving are those the evaluation will report: the RMS and largest position error, the largest tilt, the largest and the mean tension, the share of time with a tension outside $[t_{\min},t_{\max}]$, and the final position.

Without a fault the test is uneventful. Every metric changes by less than 0.03\% at every step, and so does every metric of the actuator fault and of the lower cable break, although in the latter the platform tilts by 15$^\circ$ and at least one cable is outside its bounds for 35\% of the run. The difficulty sits entirely in the peak tension after a break, because a cable that catches its load again produces a short tension peak whose height depends on where the integration step falls relative to the instant of tightening. Halving the step from 125 to 62.5~$\mu$s still changes that peak by 0.27\% after the upper cable breaks (0.35\% without noise) and by 0.15\% in the stiff case without noise; halving it from 62.5 to 31.25~$\mu$s changes it by less than 0.03\% in every scenario except one. That exception is the upper cable break without noise, at 0.12\%. In that run the verification controller, which does not know about the fault, lets the platform turn over, at least one cable is outside its bounds for half of the run, and the peak tension reaches 537~N, more than two and a half times $t_{\max}$. A run like that is a lost platform whatever the third digit of its peak tension, and we accept the 0.12\% rather than halve the step once more, which would double the cost of every run in the campaign. The plant step is therefore 62.5~$\mu$s.

The pointwise histories converge more slowly than the metrics, and we report them for that reason. Between 62.5 and 31.25~$\mu$s the position histories agree within 37~$\mu$m in the worst case, the upper cable break without noise, while the tension histories differ by up to 2.4\% of $t_{\max}$ at single samples during the snap sequences after the upper cable breaks. Pointwise tension histories during such sequences are therefore not used as evidence; the metrics above are.

\subsection{Reference solution}

Without sensor noise, the closed loop of the convergence test was solved a second time with the implicit Radau IIA method \citep{hairer1996stiff} of SciPy \citep{virtanen2020scipy}, with relative tolerance $10^{-9}$, absolute tolerance $10^{-12}$, the exact Jacobian of the compiled plant, and a restart at every inner sample. Over 6~s without a fault, RK4 at any of the steps tested agrees with it to $10^{-13}$~m. With the lower cable breaking at 4~s, RK4 at 62.5~$\mu$s stays within $5.8\times10^{-8}$~m and $2.3\times10^{-7}$~rad of the reference, and its tensions within 0.010~N, which is three orders of magnitude finer than anything the evaluation reports.

\subsection{Independent cross-check in MuJoCo}

The plant and MuJoCo \citep{todorov2012mujoco} share no code, so their agreement checks the rigid body dynamics, the cable kinematics and the unilateral cable law together. With the winches locked and the cables undamped, each cable is a spring of stiffness $EA/l_{u,i}$ that pulls only when stretched, and a spatial tendon in MuJoCo with the dead band $[0,l_{u,i}]$ on its spring length behaves in exactly that way. Both codes were started from the equilibrium at the home pose with an initial twist and run for 5~s with RK4 at 0.1~ms. With a small twist all cables stay taut and the trajectories agree to $10^{-11}$~m and $4\times10^{-11}$~rad, or $6\times10^{-9}$ of the motion. With a large twist, 0.8~m/s upward and 1~rad/s about $x$, cables go slack and tighten again repeatedly, slack in 60\,436 of the 400\,008 samples of the eight cables, and the agreement is still $8\times10^{-12}$~m and $1.4\times10^{-11}$~rad, against a motion of 59~mm and 0.09~rad. The specification asked for 1\%.

% =====================================================================
\section{Cost of a run and planning of the campaign}\label{sec:cost}

On one core of this workspace a 16~s run on T1 with the compiled plant at 62.5~$\mu$s, the sensors, the logging and the verification controller takes 11.7~s, that is 0.73~s of computing per simulated second; at 125 and 250~$\mu$s it takes 7.8 and 6.2~s. The prediction model of the estimators costs 55~$\mu$s for one RK4 step, 274~$\mu$s for a step with its Jacobians when interpreted, and 46~$\mu$s when compiled to C, and the 17 modes of the IMM together take 0.28~ms per cycle compiled. Run at 1~kHz, the IMM predictions alone thus add 0.28~s per simulated second, before the filter updates and the MPC, which do not exist yet and are timed in Phase~3.

From these figures a run of the full method should cost between 25 and 65~s of computing; the range is an estimate, not a measurement, because two of its components are still to be built. The specification asks for at least 100 runs for each of 40 single fault scenarios and the fault free case, which is 4100 runs, or 28 to 74 core hours, for one method under one condition. The proposed method under the two trajectories and the two noise levels, four baselines and three ablations make about ten blocks of that size, so the full campaign is in the order of 300 to 750 core hours. This workspace has two cores and does not keep jobs running unattended between sessions, so it cannot carry a campaign of that size; it can carry the calibration and the tests of Phase~3 and a reduced pilot campaign. For the full campaign a machine with eight or more cores that can run for several days without supervision is needed, for instance a university server. Two levers will reduce the cost once Phase~3 has timed the full method: running the IMM at 200 to 500~Hz, which the specification allows, and compiling the filter updates as the predictions are compiled here.

% =====================================================================
\section{Carried into Phase 3}\label{sec:next}

The verification controller turned the platform over after an upper cable break, which the planner of the study is designed to prevent, and this raises a point that the specification does not settle yet: when a run counts as a lost platform. We will fix a termination rule before any result is produced, for instance a tilt beyond $30^\circ$ or a drum leaving its window, so that the failed runs of the baselines end in a defined state rather than in an arbitrary flight. Phase~3 also replaces the verification controller by the estimators and the controller of the methods, tunes them on calibration runs only, and times them, which closes the cost estimate of Section~\ref{sec:cost}.

\bibliographystyle{plainnat}
\bibliography{phase2}
\end{document}
